% 第五章 神经形态计算初步探索
\chapter{神经形态计算初步探索}

磁性拓扑结构具有非易失性、低功耗和多样的动力学行为，很适合用来做神经形态计算的硬件。基于前几章对 Hopfion 稳定性和动力学调控的研究，本章从物理机制映射的角度，初步分析 Hopfion 在神经形态计算中的应用潜力。


\section{Hopfion 动力学特征与神经元功能的类比}

生物神经元的核心功能可概括为：（1）接收输入信号；（2）阈值判定后产生脉冲响应（激发）；（3）脉冲传播至下游神经元；（4）响应后进入不应期（抑制）。第四章的仿真结果揭示，Hopfion 的自旋波驱动行为与上述功能存在以下对应关系：

\begin{enumerate}
  \item \textbf{阈值效应}：幅度扫描（\cref{subsec:amplitude-scaling}）显示，自旋波幅度 $B_0$ 存在驱动阈值——$B_0 \leq \SI{0.1}{T}$ 时 Hopfion 几乎不动，$B_0 \geq \SI{0.5}{T}$ 时才产生有效运动。这与神经元的"全或无"激发特性类似。
  \item \textbf{频率选择性}：频率扫描（\cref{subsec:freq-sweep}）表明 Hopfion 对特定频率的自旋波响应强烈（如 $\SI{1100}{GHz}$），对其他频率几乎不响应（死区 $\SIrange{400}{600}{GHz}$）。这一选择性类似于神经元对特定突触输入模式的选择性响应。
  \item \textbf{方向编码}：拓扑 Hall 效应（\cref{subsec:hall-effect}）使得 srcX 和 srcZ 驱动产生完全不同的运动方向（$+z$ vs $-z$），且 Hall 角在拓扑保护下高度稳定。这为在单一器件中实现兴奋性（$+z$）和抑制性（$-z$）双通道信号编码提供了物理基础。
  \item \textbf{双向可控}：双向控制实验（\cref{subsec:bidirectional-control}）验证了通过切换自旋波频率或传播方向，可实现 Hopfion 的可逆运动，类似于神经元的激发-恢复循环。
\end{enumerate}


\section{概念器件设计}

基于上述物理映射，可以构想一种基于 Hopfion 的自旋波神经形态器件。其核心结构包括：

\begin{enumerate}
  \item \textbf{自旋波输入端}：微波天线阵列作为"突触"，通过自旋波向 Hopfion 传递输入信号。不同天线的频率、幅度和传播方向组合构成多维输入空间。
  \item \textbf{Hopfion 处理单元}：Hopfion 作为"神经元胞体"，其位置变化编码输出信号。阈值效应提供天然的非线性激活函数，频率选择性实现输入滤波。
  \item \textbf{位置读出}：通过磁阻效应（TMR/GMR）或拓扑 Hall 效应检测 Hopfion 位置变化，输出电信号。
  \item \textbf{权重调节}：自旋波幅度 $B_0$ 作为突触权重，$v \propto B_0^2$ 的标度律确保了权重与输出的确定性映射关系。
\end{enumerate}

该方案的核心优势在于：（1）Hopfion 的三维结构提供了比二维 Skyrmion 更高的存储密度和更多的动力学自由度；（2）拓扑保护确保了信号处理的鲁棒性；（3）自旋波驱动避免了电流焦耳热问题。

然而，实际实现仍面临若干挑战：Hopfion 目前仅在低温环境下被实验观测到\cite{zheng2023hopfion}，室温稳定性有待突破；竞争交换体系中的 Hopfion 特征尺寸仅 $\sim\SI{2}{nm}$，接近当前加工工艺极限；自旋波在纳米尺度传播的衰减问题也需要进一步研究。


\section{本章小结}

本章基于前几章的动力学仿真结果，从物理机制映射的角度分析了 Hopfion 在神经形态计算中的应用潜力。Hopfion 的阈值效应、频率选择性、方向编码和双向可控等动力学特征，与神经元的激发、滤波、兴奋/抑制和恢复功能有自然的物理对应。本章基于这些对应关系提出了概念器件设计方案，但从概念到实验验证还需要在材料体系、加工工艺和室温稳定性等方面取得突破。
